\section{Ethical Considerations}
\label{sec:ethical-considerations}

This work studies a formal admission substrate for capability-bounded
tool use and federated receipts. It does not report a user study,
measurement of a live population, offensive security experiment, or
deployment against production users.

\subsection{Human subjects and data}
\label{sec:ethical-considerations:subjects}

No human-subjects research and no animal-subjects research were
conducted. No IRB review or animal-care protocol was required. The
evaluation uses synthetic or replayable receipt fixtures and formal
models of admission, treaty intersection, capability attenuation, and
amendment refinement. It does not collect, process, or retain live user
data, personally identifying information, private messages, account
contents, payment records, medical records, or location traces.

\subsection{Deployment posture}
\label{sec:ethical-considerations:deployment}

The paper does not deploy an offensive tool, exploit a third-party
system, bypass an access-control policy, or probe a live service. The
construction is a verifier-side discipline for deciding when a receipt
is admitted under a constitution. Its intended use is to reduce
ambient authority and improve auditability in agent systems by making
capability scope, guard decisions, and receipt admission explicit.

\subsection{Dual-use risks}
\label{sec:ethical-considerations:dual-use}

The main dual-use risk is policy misuse by a verifier or capability
authority. A verifier could encode an exclusionary or coercive
constitution, or use the receipt log and admission predicates as a
rigid administrative control rather than as a narrowly scoped safety
mechanism. A second risk is admission-policy overclaiming: operators
could present a mechanically checked bounded model as if it proved all
runtime, cryptographic, economic, and governance properties of a live
deployment. The paper limits this risk by separating proved theorem
statements from implementation assumptions and by naming boundaries
that remain outside the formal substrate.

\subsection{Responsible release}
\label{sec:ethical-considerations:release}

The planned release is source code plus Lean artifacts, not a hosted
authority service or a production governance system. Public artifacts
will include responsible caveats: proof obligations cover the stated
model only; deployment requires independent policy review; admission
predicates should be minimized to legitimate security and audit needs;
and receipt retention should follow data-minimization and access-control
requirements appropriate to the deploying organization.
